\section{Recoil, Noncompetitive Causality, and the Architecture of Humility}

Even when the intellect grants that every creature receives its act of being continuously from God and adds nothing to him in return, a stubborn existential resistance can tighten in the chest: \emph{If all of that is true, then what am I?}

That question carries genuine human weight. If God causes my very existence at every instant, causes my mind and will to operate, and needs nothing from me, then am I anything more than an animated shadow---a puppet dangling on invisible ontological strings, a flute through which another blows every note? If everything real in me is received, does human agency dissolve into theater? Does accepting total, one-sided dependence mean accepting my own insignificance, surrendering dignity, and shrinking into servile self-contempt?

That tightening of the spirit is \emph{recoil}. It is the deep creaturely fear that God's fullness leaves no room for the creature to be genuinely real, genuinely free, and genuinely noble. In ancient theology, Augustine traced the primordial root of that recoil to the fall of the proud spirits who, ``being enamoured rather of their own power, as if they could be their own good, lapsed to this private good of their own, from that higher and beatific good which was common to all,'' forsaking ``Him who supremely is'' to turn toward ``themselves who have no such essence.''\endnote{Augustine, \emph{De Civitate Dei} XII.1 and XII.6 (Dods translation, 1871).} In modernity, the same recoil often wears a philosophical and humanist face: the conviction that human freedom and dignity can stand upright only if God is pushed back, scaled down, or eliminated altogether, because any absolute Creator would crush the autonomy of the finite self.

Yet a single hidden premise drives every form of recoil, ancient or modern. Recoil assumes that God and the creature are \emph{competitors}---that divine causality and creaturely causality operate on the same plane and divide a single, finite supply of agency between them. On that zero-sum picture, if God supplies eighty percent of an action, the creature supplies twenty percent; if God supplies one hundred percent of my being and power, then zero percent is left for me. That zero-sum picture is metaphysically false.

\subsection{Noncompetitive Causality: God Working in Every Agent}

In \emph{ST} I, q.\ 105, a.\ 5, Aquinas meets the exact objection behind recoil as his own objectors state it. If God works sufficiently in every agent, then ``it would be superfluous for the created agent to work at all'' (obj.\ 1); and since one action cannot proceed simultaneously from two sources, ``if the creature's operation is from God operating in the creature, it cannot at the same time proceed from the creature; and so no creature works at all'' (obj.\ 2).

Aquinas begins his answer by rejecting \emph{occasionalism}---the doctrine held by certain medieval Islamic and later modern thinkers that God alone truly acts while created things are mere empty occasions, ``that it is not fire that gives heat, but God in the fire, and so forth.'' He calls that position impossible on two grounds:

\begin{quote}
First, because the order of cause and effect would be taken away from created things: and this would imply lack of power in the Creator: for it is due to the power of the cause, that it bestows active power on its effect. Secondly, because the active powers which are seen to exist in things, would be bestowed on things to no purpose, if these wrought nothing through them. Indeed, all things created would seem, in a way, to be purposeless, if they lacked an operation proper to them; since the purpose of everything is its operation\dots\ We must therefore understand that God works in things in such a manner that things have their proper operation.\endnote{Aquinas, \emph{ST} I, q.\ 105, a.\ 5, corpus.}
\end{quote}

An occasionalist deity who does everything himself and leaves creatures causally hollow would display a \emph{lack of power}. The true measure of divine omnipotence is that God communicates not merely passive existence, but genuine, active causal power---making fire truly heat, oak trees truly grow, and human beings truly think, love, and decide.

The very same action proceeds wholly from God and wholly from the creature without dividing a single finite stock of agency, because primary and secondary causes belong to different orders. As Aquinas writes in Reply~2: ``One action does not proceed from two agents of the same order. But nothing hinders the same action from proceeding from a primary and a secondary agent.''\endnote{Aquinas, \emph{ST} I, q.\ 105, a.\ 5, ad 2. See also Reply~1: ``God works sufficiently in things as First Agent, but it does not follow from this that the operation of secondary agents is superfluous.''} Two men carrying a heavy timber share the same causal order: the more weight one lifts, the less the other lifts. A novelist and her characters, or a living intellect and the words it thinks, do not act on the same causal plane---and even those analogies remain pale beside the Creator, because God operates in every created agent across all three principles of action: as the ultimate \emph{end} toward which every good aims, as the \emph{first mover} applying created powers to act, and as the \emph{giver and preserver} of the very form and faculty by which the creature acts (\emph{ST} I, q.\ 105, a.\ 5, corpus).

Consequently, divine nearness and creaturely autonomy do not vary inversely; they vary directly. The closer a creature stands to God as its sustaining source, the \emph{more} real, active, and genuinely itself it becomes. God does not push my freedom aside to make room for his act; God causes my act precisely \emph{as} a free human act, upholding my intellect, my deliberation, and my willing from inside without forcing them from outside.

\subsection{Concurrence and the Boundary of Sin}

Here, however, an honest critic presses the hardest test against noncompetitive causality. If God works intimately in every created operation---including the movements of the human will---what happens when a creature sins? Either the human will truly acts on its own apart from God (breaking universal dependence), or else God causes the sinful act and becomes the author of moral evil.

Aquinas meets that dilemma directly in \emph{ST} I-II, q.\ 79, a.\ 2, building on Augustine's distinction in \emph{City of God} XII.7 between an efficient cause and a deficient cause:

\begin{quote}
Let no one, therefore, look for an efficient cause of the evil will; for it is not efficient, but deficient, as the will itself is not an effecting of something, but a defect. For defection from that which supremely is, to that which has less of being,---this is to begin to have an evil will.\endnote{Augustine, \emph{De Civitate Dei} XII.7 (Dods translation, 1871).}
\end{quote}

Evil is not a positive substance or competing essence created alongside good; it is a privation, a voluntary falling away from due order in a nature created good out of nothing. Aquinas translates that insight into the exact mechanics of divine concurrence in \emph{ST} I-II, q.\ 79, a.\ 2. Every sinful act contains two distinguishable aspects: whatever positive reality and activity belong to the act itself (speaking a word, moving a hand, choosing an apparent good), and the \emph{moral defect}---the failure of the free will to keep its act aligned with the rule of reason and divine order.

Insofar as the act is a being and an act, it is upheld by God as First Being and First Act. The defect, however, proceeds solely from the created free will falling short of due order:

\begin{quote}
Consequently this defect is not reduced to God as its cause, but to the free-will: even as the defect of limping is reduced to a crooked leg as its cause, but not to the motive power, which nevertheless causes whatever there is of movement in the limping. Accordingly God is the cause of the act of sin: and yet He is not the cause of sin, because He does not cause the act to have a defect.\endnote{Aquinas, \emph{ST} I-II, q.\ 79, a.\ 2, corpus. In Reply~2, Aquinas confirms that man is the cause of both the act and its defect, ``which defect consists in man not being subject to Whom he ought to be,'' whereas God causes the positive reality of the act without in any way causing the defect.}
\end{quote}

Creaturely freedom is therefore neither a rival mini-aseity standing outside God's conserving power nor a sham. It is so genuinely real that a rational creature can actually deflect its own act away from the order of truth and good, bearing full moral responsibility for the wound it introduces, even while the very power to act remains every second a gift from God. Recoil is not a sign of superior independence; it is the gesture by which a finite will misuses borrowed strength to pretend it is unborrowed.

\subsection{Humility and Magnanimity: Truth Without Self-Contempt}

Even when noncompetitive causality secures creaturely agency in principle, the very word \emph{humility} often triggers quiet dread. To modern sensibilities shaped by Enlightenment pride and Nietzschean suspicion, humility sounds like self-erasure---cringing, hiding one's capacities, calling oneself worthless, and retreating from great endeavors lest one seem presumptuous. If accepting creaturehood means cultivating a small, frightened soul, recoil looks almost heroic by comparison.

Aquinas dismantles that caricature by setting \emph{humility} (\emph{ST} II-II, q.\ 161, aa.\ 1--3) directly alongside \emph{magnanimity} (\emph{ST} II-II, q.\ 129, aa.\ 1, 3, ad 4). Far from opposing each other, these two virtues form a single, balanced architecture of creaturely truth.

For Aquinas, humility is neither an assault on truth nor theatrical self-deprecation---which he dismisses, quoting Augustine, as ``false humility'' and ``grievous pride'' (\emph{ST} II-II, q.\ 161, a.\ 1, ad 2). Still less does humility consist in lowering oneself below one's God-given dignity into beastliness (\emph{ST} II-II, q.\ 161, a.\ 1, ad 1). Properly speaking, humility resides in the appetite, tempering and restraining the mind ``lest it tend to high things immoderately,'' guided by truthful self-knowledge---``knowledge of one's own deficiency belongs to humility, as a rule guiding the appetite'' (\emph{ST} II-II, q.\ 161, aa.\ 1--2). Its deepest root is neither social timidity nor low self-esteem, but reverent truth before God: ``humility, considered as a special virtue, regards chiefly the subjection of man to God'' (\emph{ST} II-II, q.\ 161, a.\ 1, ad 5; a.\ 2, ad 3).

In \emph{ST} II-II, q.\ 161, a.\ 3, Aquinas separates what belongs to God from what belongs to the creature:

\begin{quote}
We may consider two things in man, namely that which is God's, and that which is man's. Whatever pertains to defect is man's: but whatever pertains to man's welfare and perfection is God's\dots\ For those who have a share of God's gifts know that they have them, according to 1~Cor.\ 2:12: ``That we may know the things that are given us from God.''\endnote{Aquinas, \emph{ST} II-II, q.\ 161, a.\ 3, corpus. Compare St.\ Paul's searching question in 1~Cor.\ 4:7 (Douay-Rheims): ``Or what hast thou that thou hast not received, and if thou hast received, why dost thou glory, as if thou hadst not received it?''}
\end{quote}

Humility never asks a person to lie about the intellect, courage, beauty, skill, or grace he genuinely possesses. To pretend that God's real gifts in you are nothing is not humility; it is ingratitude dressed in pious rags. Humility simply refuses to steal the credit for them---refuses to claim as unborrowed what was received.

Right beside humility stands \emph{magnanimity}---literally ``greatness of soul,'' which ``by its very name denotes stretching forth of the mind to great things'' (\emph{ST} II-II, q.\ 129, a.\ 1, corpus). Where humility curbs presumptuous overreaching based on self-reliance, magnanimity fortifies the soul against despair and pusillanimity, urging it toward noble, arduous works worthy of honor according to right reason (\emph{ST} II-II, q.\ 161, a.\ 1, corpus; q.\ 129, a.\ 1, ad 3). When an objector claims that humility and magnanimity contradict each other---since the magnanimous person deems himself worthy of great things while the humble person thinks little of himself---Aquinas reconciles the two virtues directly:

\begin{quote}
There is in man something great which he possesses through the gift of God; and something defective which accrues to him through the weakness of nature. Accordingly magnanimity makes a man deem himself worthy of great things in consideration of the gifts he holds from God: thus if his soul is endowed with great virtue, magnanimity makes him tend to perfect works of virtue; and the same is to be said of the use of any other good, such as science or external fortune. On the other hand, humility makes a man think little of himself in consideration of his own deficiency\dots\ It is therefore evident that magnanimity and humility are not contrary to one another, although they seem to tend in contrary directions, because they proceed according to different considerations.\endnote{Aquinas, \emph{ST} II-II, q.\ 129, a.\ 3, ad 4. See also \emph{ST} II-II, q.\ 161, a.\ 2, ad 2: ``It is contrary to humility to aim at greater things through confiding in one's own powers: but to aim at greater things through confidence in God's help, is not contrary to humility.''}
\end{quote}

Recoil assumed a tragic bargain: either be dependent and small, or be independent and great. Aquinas exposes that bargain as an illusion. Humility anchors the truth of non-aseity; magnanimity honors the reality of received gifts. Together they let the creature stand upright without pretending to have made itself.

\subsection{Clay in the Hands of the Maker}

Nearly eleven centuries before Aquinas wrote those articles, St.\ Irenaeus addressed the impatient human heart that resents having been created finite rather than already divine, in \emph{Adversus Haereses} IV.39.2:

\begin{quote}
For thou dost not make God, but God thee. If, then, thou art God's workmanship, await the hand of thy Maker which creates everything in due time; in due time as far as thou art concerned, whose creation is being carried out. Offer to Him thy heart in a soft and tractable state, and preserve the form in which the Creator has fashioned thee, having moisture in thyself, lest, by becoming hardened, thou lose the impressions of His fingers\dots\ But if thou, being obstinately hardened, dost reject the operation of His skill, and show thyself ungrateful towards Him, because thou wert created a [mere] man, by becoming thus ungrateful to God, thou hast at once lost both His workmanship and life. For creation is an attribute of the goodness of God but to be created is that of human nature.\endnote{Irenaeus, \emph{Adversus Haereses} IV.39.2 (ANF, vol.\ 1, Roberts and Rambaut translation).}
\end{quote}
